\scanpage{5}
degrés positifs. On voit que si $i$ est effectif, il en est de $=$ de ses conjugués pr éléments de $G(\mathbb{C})$, et de ses conjugués pr autom. de $\mathbb{C}$, en particulier de $\bar i$. D'autre part, pour $i$ donné, la sous-catégorie \inserted{pleine} $P(i)$ de $M$ formée des $E$ tels que $E_{\mathbb{C}}$ soit à $i$-\unsure{graduations} positives est \inserted{str. pleine,} stable pr \deleted{$\unsure{\otimes}$} sommes, pr \deleted{\unsure{les}} sous-objets, pr $\otimes$, contient \inserted{$0$, et} $\Lambda(-1)$ si on suppose que \deleted{$\varepsilon i(\lambda)=\lambda^\nu$ avec $\nu$} $\deleted{\langle i,\varepsilon\rangle}\inserted{\varepsilon(i)}\geq0$, et vérifie la condition que
\[
\forall E\in\mathrm{Ob}\,M,\quad \exists n\text{ tel que }E(-n)\in\mathrm{Ob}\,P(i),
\]
si on suppose $\boxed{\deleted{\langle i,\varepsilon\rangle}\inserted{\varepsilon(i)}>0}$ [N.B. on désigne par $\langle i,\varepsilon\rangle$ l'entier $\nu$ tel que $\varepsilon(i(\lambda))=\lambda^\nu$].

\deleted{Si \unsure{on prend} \unsure{conjugué par} les conjugués $\operatorname{int}(g)\circ i$ et les conjugués $\bar i$ $(g\in G(\mathbb{C}),\,\unsure{i\in})$}
\inserted{Prenant des intersections} de sous-catégories $P(i)$ on trouve des $M^+$ satisfaisant aux axiomes considérés, sauf peut être \unsure{au dernier} sur la graduation.
\inserted{Par intersection de tels \unsure{foncteurs}, itou\ldots}
Les $i$ pour lesquels, $M^+$ étant donné,

\mspage{4}
